\subsection{十四类章节的写法}

\begin{longtable}{p{1.8cm}p{4.2cm}p{4.0cm}p{4.0cm}}
\caption{章节功能、段落套路与质量门}\label{tab:sections}\\
\toprule
章节 & 必答问题与段落结构 & 推荐表达及图表 & 质量门与常见错误\\
\midrule
\endfirsthead
\toprule
章节 & 必答问题与段落结构 & 推荐表达及图表 & 质量门与常见错误\\
\midrule
\endhead
摘要 & 逐问写对象、模型、关键结果和检验；最后给关键词。 & “针对……，建立……；求得……；经……检验……”；通常不放图。 & 每问均出现，数字与正文一致；不堆算法名，不写“效果良好”。\\
问题重述 & 重组给定量、输出、硬约束、附件和问间依赖。 & “给定……，要求在……约束下确定……”；复杂几何可重绘示意图。 & 不大段照抄，不漏单位、时段或提交文件，不在前文引入模型。\\
问题分析 & 每问说明核心矛盾、输入输出、基线、改进和数据流。 & “由于……，故先……，再……”；节末可放总流程图。 & 必须解释选择理由和近邻方法边界，不能只列算法。\\
模型假设 & 一条一义写对象、范围、简化理由和方程后果。 & “在……尺度内忽略……，因此……”；必要时放情景表。 & 题面事实不是假设；不可写“数据真实可靠”；后文可追溯。\\
符号说明 & 定义状态量、决策量、参数、集合、索引、单位和域。 & 三线表分组；局部临时量在首次出现处定义。 & 同符号不跨问改义；公式、代码和附件字段一致。\\
模型建立 & 依次给依据、变量、数学表达、约束、初边值和参数。 & “令……；由……可得……”；结构图靠近定义。 & 公式前有来源、后有解释；约束方向、量纲和变量域正确。\\
模型求解 & 写预处理、初始化、求解器/迭代、参数、停止和输出。 & 按执行顺序写；伪代码和流程图不替代参数表。 & 第三方可复算；不以“调用 MATLAB/Python”结束。\\
结果分析 & 读关键值，比较基线，解释原因，说明约束和决策含义。 & “在……条件下，……由……变为……，表明……”；图表后立即解释。 & 不逐格念表；有单位、有效数字和不确定范围；每问均回应。\\
模型检验 & 明确检验对象、方法、阈值、独立数据、指标和裁决。 & 回代、守恒、残差、网格收敛、交叉验证或消融。 & 不能用训练拟合代替外部能力；失败项也要报告。\\
灵敏度分析 & 定义扰动参数、范围、采样、重复求解和响应指标。 & “当……位于……时结论不变；超过……后策略切换”；放曲线/情景表。 & 扰动后必须重算；区分数值稳定和决策稳定。\\
模型评价 & 将优势绑定机制，将不足绑定假设、数据或算法。 & “由于显式保留……，模型能够……；受……限制……” & 不写任何模型都适用的“简单准确、适用面广”。\\
改进方案 & 对每项局限给新数据/变量/约束、实现方式和验证协议。 & 按“缺陷—改法—验证”对应；可用短表。 & 增加复杂度不是天然改进；方案必须可执行和可检验。\\
参考文献 & 只列正文实际使用的理论、数据、软件与标准来源。 & 统一著录；定义和算法优先原始论文、教材、标准或官方资料。 & 正文引用与条目双向对应；网页含机构、日期和访问日期。\\
附录 & 给入口脚本、环境、参数、随机种子、数据字典和补充结果。 & 文件清单、目录树、可复制代码；主文明确引用。 & 关键结果可由附录复现；不放截图代码，不泄露个人信息。\\
\bottomrule
\end{longtable}

